% !PRINTED 532
\weberpage{532}

Aus (27) leitet man eine einfachere Formel her für die Funktion
\[
f(\omega) = e^{-\frac{\pi i}{24}}\,\frac{\eta\!\left(\frac{\omega+1}{2}\right)}{\eta(\omega)} = e^{\frac{\pi i}{24}}\,\frac{\eta\!\left(\frac{\omega-1}{2}\right)}{\eta(\omega)}.
\]

Ersetzt man nämlich
\[
\begin{aligned}
A &\ \text{durch}\ 2A,\\
B &\ \text{\textquotedblbase}\ \ A+B,\\
C &\ \text{\textquotedblbase}\ \ \tfrac12(A+2B+C),
\end{aligned}
\]
so bleibt $m=AC-B^2$ ungeändert und $\omega_1,\omega_2$ gehen über in
\[
\omega_1' = \frac{\omega_1+1}{2},\quad \omega_2' = \frac{\omega_2-1}{2}.
\]

Setzt man also
\weq{29}{\begin{aligned}
\psi &= (A,\,2B,\,C),\\
\psi' &= [2A,\,2(A+B),\,\tfrac12(A+2B+C)],
\end{aligned}}
so ergibt sich aus (27)
\weq{30}{\mathop{\mathrm{Lim}}\limits_{s=1}\left(\sum^{x,y}\frac{1}{\psi^s} - \sum^{x,y}\frac{1}{\psi'^s}\right) = \frac{2\pi}{\sqrt m}\log\frac{f(\omega_1)f(\omega_2)}{\sqrt 2}.}

Wir nehmen jetzt nicht nur $A, 2B, C$, sondern auch $A, B, C$ als ganze Zahlen ohne gemeinschaftlichen Teiler und $A$ ungerade an. Dann ist $\psi$ eine primitive Form der Diskriminante $-4m$, und wir betrachten zunächst die beiden Fälle
\[
m\equiv 1,\quad m\equiv 2\ (\mathrm{mod}\ 4).
\]

Nehmen wir, was keine weitere Beschränkung ist, $B\equiv m+1\ (\mathrm{mod}\ 2)$ an, d. h. $B$ gerade oder ungerade, je nachdem $m$ ungerade oder gerade ist, so ergibt sich aus
\[
m = AC-B^2\text{:}
\]
\[
\begin{aligned}
A-C&\equiv 0\ (\mathrm{mod}\ 4) \quad m\ \text{ungerade},\\
A+C&\equiv 0\ \text{\textquotedblbase} \qquad\ \ m\ \text{gerade},
\end{aligned}
\]
und die Form $\psi'$ ist primitiv von der Diskriminante $-4m$. Es ist nach der Komposition der quadratischen Formen:
\[
\psi' = \psi_0\psi,
\]
wenn
\[
\begin{aligned}
\psi_0 &= \left(2,\,2,\,\tfrac{m+1}{2}\right) & (m\ \text{ungerade}),\\
&= \left(2,\,0,\,\tfrac{m}{2}\right) & (m\ \text{gerade}),
\end{aligned}
\]
%CONT>>
